% ID: FIS-OTT-SPE-001
% Titolo: Immagine ingrandita con specchio concavo
% Disciplina: Fisica
% Area: Ottica
% Argomento: Specchi
% Sottoargomento: Specchio concavo e ingrandimento
% Classe: Seconda liceo scientifico
% Difficolta: 3
% Tipo: problema_numerico
% Risultato: 8 cm davanti allo specchio
% Tempo_stimato: 8 min
% Competenze: uso dell'equazione degli specchi, interpretazione dell'ingrandimento, modellizzazione
% Prerequisiti: specchi sferici, distanza focale, ingrandimento
% Tag: fisica, ottica, specchi, specchio_concavo, ingrandimento
% Fonte: verifica_fisica_ottica_anonimizzata
% Autore: Riccardo Carli
% Licenza: CC-BY-SA-4.0

\begin{esercizio}
Utilizzando uno specchio concavo di distanza focale \(12\,\mathrm{cm}\), si
vuole ottenere un'immagine diritta di uno spillo che abbia dimensioni triple
rispetto a quella reale.

In quale punto dell'asse ottico deve essere collocato lo spillo?
\end{esercizio}

\begin{soluzione}
Per uno specchio concavo vale l'equazione:
\[
\frac{1}{f}=\frac{1}{p}+\frac{1}{q},
\]
dove \(p\) e la distanza dell'oggetto e \(q\) quella dell'immagine.

L'immagine richiesta e diritta e tre volte piu grande, quindi
l'ingrandimento e:
\[
M=+3.
\]
Per gli specchi:
\[
M=-\frac{q}{p}.
\]
Allora:
\[
3=-\frac{q}{p}
\qquad\Rightarrow\qquad
q=-3p.
\]
Il segno negativo indica un'immagine virtuale, dietro lo specchio.

Sostituendo nell'equazione degli specchi:
\[
\frac{1}{f}=\frac{1}{p}+\frac{1}{-3p}
=\frac{2}{3p}.
\]
Quindi:
\[
p=\frac{2}{3}f
=\frac{2}{3}\cdot 12\,\mathrm{cm}
=8\,\mathrm{cm}.
\]

Lo spillo deve essere collocato a \(8\,\mathrm{cm}\) davanti allo specchio,
cioe tra il vertice dello specchio e il fuoco.
\end{soluzione}
